% ch05-hawkes — exponential Hawkes: excitation, inhibition, and covariates

\section{Model and assumptions}\label{ch05:sec:model}

The Poisson models of \cref{ch:poisson} treat deal flow as exogenously
modulated: nothing that happens in the market changes the future of the
market. Their residual diagnostics say otherwise --- clustering beyond what
covariates explain is the recurring finding of the rehearsal campaigns
(\cref{app:results}) --- and the Hawkes process \citep{hawkes1971} is the
minimal extension in which events beget events. This chapter develops the
exponential-kernel Hawkes layer end to end: the branching representation
that organizes everything (\cref{ch05:sec:cluster}), the likelihood and the
recursion that makes it computable (\cref{ch05:sec:likelihood}), the EM
algorithm on the latent genealogy (\cref{ch05:sec:em}), simulation
(\cref{ch05:sec:simulation}), and then the section this chapter exists for:
what to do about \emph{inhibition}, which a nonnegative-kernel Hawkes
process provably cannot express, and which our market exhibits --- a firm
that has just raised is temporarily \emph{less} likely to raise again
(\cref{ch05:sec:inhibition}). The chapter closes with the
frailty-versus-contagion identifiability problem
(\cref{ch05:sec:frailty}) and the defect declaration required by D2.1
(\cref{ch05:sec:defects}).

\begin{assumption}[Exponential Hawkes ground process]\label{asm:ch05:model}
On the filtered space of \cref{def:ch02:market}, the univariate model for a
single stream is
\begin{equation}\label{eq:ch05:uni}
\lam(t) \;=\; \mub(t) \;+\; \sum_{t_m < t} \kernel(t - t_m),
\qquad
\kernel(u) \;=\; \branch\decay\, e^{-\decay u},
\quad \branch \ge 0,\ \decay > 0,
\end{equation}
and the multivariate, sector-indexed model is, for $s \in \sectors$,
\begin{equation}\label{eq:ch05:multi}
\lam_s(t) \;=\; \mub_s(t) \;+\; \sum_{r \in \sectors}\;
\sum_{\substack{t_m < t \\ s_m = r}}
\branch_{sr}\,\decay_{sr}\, e^{-\decay_{sr}(t - t_m)},
\qquad
\mub_s(t) \;=\; \exp\!\big(\bcoef_s^{\top} \xcov_s(t)\big),
\end{equation}
where the covariate process $\xcov_s(t)$ (macro series, sector activity
summaries, seasonal terms) is predictable and locally bounded, constructed
under the point-in-time protocol \cref{prot:ch02:pit}. The
\emph{branching matrix} is
\begin{equation}\label{eq:ch05:gmat}
\Gmat_{sr} \;=\; \int_0^{\infty} \branch_{sr}\decay_{sr} e^{-\decay_{sr}u}
\dd u \;=\; \branch_{sr},
\end{equation}
and we impose \emph{subcriticality}: $\specrad{\Gmat} < 1$.
\end{assumption}

Two virtues of the $(\branch,\decay)$ parameterization justify committing
to it over the equally common $\alpha e^{-\beta t}$ form. First,
$\branch = \int_0^\infty \kernel$ is a \emph{unitless branching mass}: by
\cref{ch05:sec:cluster} it is the expected number of direct offspring per
event, invariant to the time unit, while $\decay$ is a pure timescale with
half-life $\log 2/\decay$. Second, the two \emph{decouple in aggregation
analysis}: binning event times at width $\horizon$ (defect (C4)) destroys
information about $\decay$ once $\horizon\decay$ is large, but the mass
$\branch$ remains estimable from count balance --- the
$\branch$--$\decay$ ridge quantified in \cref{ch:gof} and the standing
verdict ``$\branch$ survives aggregation, $\decay$ does not''
(\cref{app:results}). A parameterization in which the kernel height
$\alpha = \branch\decay$ is a free parameter entangles mass and timescale
in exactly the direction the data cannot resolve; keeping
$(\branch,\decay)$ explicit keeps the ill-conditioned direction visible.
(The event-time fitters currently store $\alpha=\branch\decay$ and $\beta
= \decay$ internally; the mapping is recorded in the code hooks.)

Subcriticality is the standing stability constraint. For constant
baselines it is exactly the condition for a stationary version to exist
\citep{hawkes1971,bremaud1996}; for the covariate-driven baseline of
\cref{eq:ch05:multi} the process is defined conditionally on the covariate
paths and is non-explosive on finite windows regardless
(\cref{rem:ch05:nonexplosion}), but $\specrad{\Gmat}<1$ still governs
cluster sizes, forecast amplification, and the sanity of any simulation.
Because $\Gmat \ge 0$, Perron--Frobenius gives $\specrad{\Gmat} \le \max_s
\sum_r \Gmat_{sr}$, so a row-sum bound is a convex sufficient condition
--- this is precisely the constraint the repository's discrete-face fitter
enforces with $\rhomax = 0.95$ (\modrepo{sector\_stability.py}), and the
continuous-time layer inherits the same certificate. Multivariate Hawkes
models with spectral-radius stability conditions are standard in financial
applications \citep{embrechts2011,bacry2015}.

\section{The cluster representation}\label{ch05:sec:cluster}

Everything tractable about the linear Hawkes process follows from one
theorem: it is secretly a Poisson cluster process
\citep{hawkesoakes1974}.

\begin{theorem}[Hawkes--Oakes cluster representation]\label{thm:ch05:cluster}
Let $\mub_s(t) \ge 0$ be locally integrable baselines and $\kernel_{sr}
\ge 0$ kernels with $\Gmat_{sr} = \int \kernel_{sr} < \infty$. Construct a
marked point process as follows: \emph{immigrants} of type $s$ arrive as
an inhomogeneous Poisson process with rate $\mub_s(t)$, independently
across types; recursively, each event of type $r$ at time $t_m$ (immigrant
or offspring) generates \emph{direct offspring} of type $s$ on
$(t_m,\infty)$ as an inhomogeneous Poisson process with rate
$\kernel_{sr}(t - t_m)$, independently of everything else; the process is
the superposition of all generations. Then the superposition is a Hawkes
process with intensity \cref{eq:ch05:multi}. Conversely, every Hawkes
process with these parameters is equal in law to this cluster process.
\end{theorem}

\begin{proof}[Proof of the construction direction]
Let $\{\mathcal{G}_t\}$ be the filtration generated by the labeled
construction (all events with their genealogy). Fix $t$ and consider
arrivals of type $s$ in $(t, t+\dd t]$. They are of two kinds. New
immigrants form a Poisson process with deterministic rate $\mub_s$,
independent of $\mathcal{G}_t$, hence contribute $\mathcal{G}_t$-intensity
$\mub_s(t)$. For each event $t_m < t$ of type $r$ already born, its
type-$s$ offspring process is inhomogeneous Poisson with rate
$\kernel_{sr}(\cdot - t_m)$; by independent increments of the Poisson
process, its future beyond $t$ is conditionally independent of
$\mathcal{G}_t$ with intensity $\kernel_{sr}(t - t_m)$. Events born after
$t$ contribute nothing, since the kernels are supported on positive ages.
The streams are independent by construction, and the intensity of a
superposition of conditionally independent streams is the sum of the
intensities (additivity of compensators), so the type-$s$ superposed
counting process has $\mathcal{G}_t$-intensity $\mub_s(t) + \sum_r
\sum_{t_m<t,\,s_m=r} \kernel_{sr}(t - t_m)$. Finally, this expression is a
functional of the unlabeled event times alone, hence measurable with
respect to the internal filtration $\Ft \subseteq \mathcal{G}_t$; by the
innovation (projection) theorem the $\Ft$-intensity is the optional
projection of the $\mathcal{G}_t$-intensity, which here equals itself
\citep[Ch.~7]{daley2003}. So the unlabeled superposition has the Hawkes
intensity \cref{eq:ch05:multi}. The converse (existence and uniqueness in
law of the cluster decomposition for a given Hawkes law) is
\citet{hawkesoakes1974}; we use only the construction direction, which is
what simulation and EM consume.
\end{proof}

\begin{remark}[Non-explosion does not need subcriticality]\label{rem:ch05:nonexplosion}
On a finite window $[0,\Tend]$ the mean intensity $\xi_s(t) =
\E[\lam_s(t)]$ satisfies the Volterra equation $\xi = \mub + \kernel *
\xi$ componentwise; since the exponential kernels are bounded by
$\branch_{sr}\decay_{sr}$, Gr\"onwall's inequality gives $\xi_s(t) \le
\bar\mub\, e^{c t} < \infty$ for constants depending only on the
parameters. The process is therefore well defined and a.s.\ finite on any
finite window for \emph{any} $\branch$; what $\specrad{\Gmat}<1$ buys is
stationarity, finite cluster sizes in expectation, and bounded long-run
amplification --- the subject of the corollaries below. This distinction
matters later: for the \emph{nonlinear} models of
\cref{ch05:sec:inhibition} even finite-window good behavior is no longer
automatic.
\end{remark}

\begin{corollary}[Subcriticality and cluster size]\label{cor:ch05:subcritical}
In the cluster construction, the offspring cascade of a single type-$r$
event is a multitype Galton--Watson process with mean-offspring matrix
$\Gmat$. The expected total cluster size (all generations, all types,
including the ancestor) is finite iff $\specrad{\Gmat} < 1$, in which case
the expected number of type-$s$ events in the cluster of a type-$r$
ancestor is $\big[(I - \Gmat)^{-1}\big]_{sr}$. In the univariate case the
expected total cluster size is $1/(1-\branch)$.
\end{corollary}

\begin{proof}
The number of direct type-$s$ offspring of a type-$r$ event is Poisson
with mean $\int \kernel_{sr} = \Gmat_{sr}$, independently across events,
so generation sizes form a multitype branching process: the expected
type-composition vector of generation $n$ descending from one type-$r$
ancestor is $\Gmat^{n} e_r$ (induction on $n$, using linearity of
expectation and independence of offspring processes). The expected cluster
composition is $\sum_{n \ge 0} \Gmat^{n} e_r$. For a nonnegative matrix
the Neumann series $\sum_n \Gmat^n$ converges iff $\specrad{\Gmat} < 1$
(Gelfand's formula), with limit $(I-\Gmat)^{-1}$; if $\specrad{\Gmat} \ge
1$ the series diverges because $\mathbf{1}^\top \Gmat^n e_r \not\to 0$
along the Perron direction. Univariate: $\sum_n \branch^n =
1/(1-\branch)$.
\end{proof}

\begin{corollary}[Stationary rate and endogenous share]\label{cor:ch05:endogenous}
Let the baselines be constant, $\mub_s(t) \equiv \bar\mub_s$, and
$\specrad{\Gmat} < 1$. The stationary mean rate vector is
\begin{equation}\label{eq:ch05:statrate}
\bar\lam \;=\; (I - \Gmat)^{-1} \bar\mub,
\qquad\text{univariate: } \bar\lam = \frac{\bar\mub}{1 - \branch},
\end{equation}
and the stationary fraction of events that are offspring rather than
immigrants (the \emph{endogenous share}) is $1 - \mathbf{1}^\top\bar\mub /
\mathbf{1}^\top\bar\lam$; in the univariate case it equals $\branch$.
\end{corollary}

\begin{proof}
Taking expectations in \cref{eq:ch05:multi} under stationarity,
$\bar\lam_s = \bar\mub_s + \sum_r \Gmat_{sr}\bar\lam_r$ because
$\E \sum_{t_m<t, s_m=r} \kernel_{sr}(t-t_m) = \int_0^\infty
\kernel_{sr}(u)\, \bar\lam_r \dd u$ by Campbell's formula; solving gives
\cref{eq:ch05:statrate}. Immigrants arrive at rate $\bar\mub_s$ by
construction, so the immigrant share of events is
$\mathbf{1}^\top\bar\mub/\mathbf{1}^\top\bar\lam$ and the endogenous share
is its complement; univariate, $1 - \bar\mub/\bar\lam = 1 - (1-\branch) =
\branch$. Equivalently: each immigrant's cluster has expected size
$1/(1-\branch)$ (\cref{cor:ch05:subcritical}), of which exactly one event
is the immigrant.
\end{proof}

\begin{remark}[A common accounting error]\label{rem:ch05:accounting}
It is tempting to reason ``rate $\bar\mub$ of immigrants, each with
$\branch$ children, hence total rate $\bar\mub(1+\branch)$ and endogenous
share $\branch/(1+\branch)$.'' This counts only the first generation:
children beget children, the geometric series compounds to
$1/(1-\branch)$, and the correct endogenous share is $\branch$, not
$\branch/(1+\branch)$. At $\branch = 0.5$ the difference is $50\%$ versus
$33\%$ of deal flow attributed to contagion --- a headline-sized error.
All endogeneity reporting in this repository uses
\cref{cor:ch05:endogenous}, and (better) the posterior parentage
diagnostic of \cref{ch05:sec:em}, subject to the upper-bound discipline of
\cref{ch05:sec:frailty}.
\end{remark}

\section{Likelihood and the Ogata recursion}\label{ch05:sec:likelihood}

The Jacod likelihood \cref{eq:ch02:jacod} specializes to
\cref{eq:ch05:uni} as
\begin{equation}\label{eq:ch05:loglik}
\loglik(\bcoef, \branch, \decay)
\;=\;
\sum_{m=1}^{n} \log\!\Big(\mub(t_m) + \branch\decay\, A(m)\Big)
\;-\; \int_0^{\Tend} \mub(u)\dd u
\;-\; \branch \sum_{m=1}^{n} \Big(1 - e^{-\decay(\Tend - t_m)}\Big),
\end{equation}
where
\begin{equation}\label{eq:ch05:Adef}
A(m) \;\defeq\; \sum_{l < m} e^{-\decay\,(t_m - t_l)}
\end{equation}
is the decayed event count seen by event $m$. The compensator term follows
by exact integration: $\int_{t_m}^{\Tend} \branch\decay
e^{-\decay(u-t_m)}\dd u = \branch(1 - e^{-\decay(\Tend-t_m)})$, summed
over events; the baseline integral is available in closed form for the
piecewise-constant covariates used in the repository
(\modrepo{eventtime/covariates.py}). Computed naively, the $A(m)$ cost
$O(n^2)$; the exponential kernel --- alone among common kernels --- makes
them recursive \citep{ozaki1979}.

\begin{proposition}[Ogata recursion]\label{prop:ch05:recursion}
With $A(1) = 0$, the quantities \cref{eq:ch05:Adef} satisfy
\begin{equation}\label{eq:ch05:recursion}
A(m) \;=\; e^{-\decay\,(t_m - t_{m-1})}\,\big(1 + A(m-1)\big),
\qquad m \ge 2,
\end{equation}
so \cref{eq:ch05:loglik} is computable in $O(n)$ after sorting.
\end{proposition}

\begin{proof}
By induction on $m$. Base case $m=2$: $A(2) = e^{-\decay(t_2-t_1)} =
e^{-\decay(t_2-t_1)}(1 + A(1))$ since $A(1)=0$. Inductive step: assume
$A(m-1)$ equals \cref{eq:ch05:Adef} at index $m-1$. Splitting off the
$l=m-1$ term and factoring the common decay over $[t_{m-1}, t_m]$,
\[
A(m)
= \sum_{l \le m-1} e^{-\decay(t_m - t_l)}
= e^{-\decay(t_m - t_{m-1})}\Big(e^{0} + \sum_{l < m-1}
e^{-\decay(t_{m-1} - t_l)}\Big)
= e^{-\decay(t_m - t_{m-1})}\big(1 + A(m-1)\big),
\]
which is \cref{eq:ch05:recursion}. The recursion works because the
exponential kernel satisfies $\kernel(u+v) = \kernel(u)e^{-\decay v}$:
the whole past compresses into one number that decays deterministically
between events. This is the Markov-embedding property; it fails for
power-law kernels, which admit no finite recursion, and it is why this
document fixes the exponential family as the parametric default.
\end{proof}

In the multivariate model the same recursion runs per source--target pair:
$A_{sr}(m) = e^{-\decay_{sr}(t_m - t_{m-1})}\big(\ind{s_{m-1}=r} +
A_{sr}(m-1)\big)$ on the merged event stream, giving $O(n\abs{\sectors}^2)$
per likelihood pass; \modrepo{eventtime/likelihood.py} implements exactly
this, in a block-cumsum form that is numerically stable over long
horizons, together with analytic gradients. Asymptotic normality of the
MLE in the stationary ergodic case is \citet{ogata1978}.

\paragraph{Positivity transforms.} The fitters optimize unconstrained by
reparameterizing $\branch = e^{a}$ (in code, $\alpha = e^{a}$ with
$\alpha=\branch\decay$), so gradient methods never propose negative
branching and standard errors return via the delta method
(\modrepo{eventtime/estimate.py}). Two consequences deserve note. First,
the transform is a \emph{modeling} statement, not a numerical trick: it
hard-codes the nonnegative-kernel assumption whose consequences
\cref{ch05:sec:inhibition} confronts --- a very negative $a$ is
near-zero excitation, never inhibition. Second, the boundary $\branch = 0$
is pushed to $a \to -\infty$, so ``no excitation'' is a flat direction of
the transformed likelihood, and testing $H_0\colon \branch = 0$ is a
boundary problem whose LRT null is the mixture
$\tfrac12\chi^2_0 + \tfrac12\chi^2_1$ \citep{selfliang1987}; the testing
protocol lives in \cref{ch:gof}. When the number of sectors makes the
$\abs{\sectors}^2$ branching entries data-hungry, the $\ell_1$-penalized
fitter with exact zeros (\modrepo{eventtime/lasso.py}) implements the
penalized multivariate excitation recovery of \citet{hansen2015}.

\begin{warning}[The likelihood is not concave in
$(\branch,\decay)$]\label{warn:ch05:nonconcave}
For \emph{fixed} $\decay$ and a linearly parameterized baseline,
\cref{eq:ch05:loglik} is concave: each event term is the log of an affine
function of the parameters and the compensator is linear. Jointly in
$(\branch,\decay)$ it is not --- the event terms involve
$\branch\decay e^{-\decay \Delta}$, and the likelihood surface develops
the $\branch$--$\decay$ ridge analyzed in \cref{ch:gof}: data observed at
resolution coarser than the kernel half-life constrain roughly the product
of mass and effective within-window decay, not each factor. With the
exp-link baseline of \cref{eq:ch05:multi} even the fixed-$\decay$
problem loses global concavity (the event terms
$\log(e^{\bcoef^\top\xcov}+c)$ are convex in $\bcoef$), retaining
concavity only in the excitation block given the baseline. Practice
therefore fixes a grid of decays, maximizes at each grid point with
multistart, and reports the \emph{profile} likelihood over $\decay$ ---
flagging ridge flatness rather than hiding it. The EM algorithm of the
next section is the complementary tool: it restores block-concavity
through augmentation and yields the diagnostics no direct MLE provides.
\end{warning}

\section{EM on the latent genealogy}\label{ch05:sec:em}

\Cref{thm:ch05:cluster} says every event has a latent \emph{parent}:
either the baseline (immigrant) or an earlier event. Making the genealogy
explicit turns Hawkes estimation into a missing-data problem with an exact
E-step and closed-form M-steps --- the cascades-of-Poisson-processes
formulation of \citet{simmajordan2010}, identical in structure to the ETAS
EM of \citet{veenschoenberg2008}. We derive it in the univariate model;
the multivariate bookkeeping is indicated at the end.

Let $u_m \in \{0, 1, \dots, m-1\}$ denote the parent of event $m$ ($0$ =
immigrant, $l \ge 1$ = event $t_l$).

\begin{lemma}[Complete-data likelihood]\label{lem:ch05:completedata}
Given the genealogy $u = (u_1,\dots,u_n)$, the complete-data
log-likelihood on $[0,\Tend]$ is
\begin{equation}\label{eq:ch05:completedata}
\loglik_c(\bcoef,\branch,\decay; u)
=
\sum_{m: u_m = 0} \log \mub(t_m)
- \int_0^{\Tend}\!\! \mub
+ \sum_{m: u_m = l \ge 1}
\big[\log\branch + \log\decay - \decay(t_m - t_l)\big]
- \branch \sum_{l=1}^{n} \big(1 - e^{-\decay(\Tend - t_l)}\big).
\end{equation}
\end{lemma}

\begin{proof}
Under the cluster construction the labeled data decompose into
independent inhomogeneous Poisson processes: the immigrant process with
rate $\mub$ on $[0,\Tend]$, contributing $\sum_{u_m=0}\log\mub(t_m) -
\int_0^{\Tend}\mub$; and, for each event $l$, its offspring process with
rate $\kernel(\cdot - t_l)$ observed on $(t_l, \Tend]$, contributing
$\sum_{m: u_m=l}\log\kernel(t_m-t_l) - \int_0^{\Tend - t_l}\kernel(v)\dd v$.
Summing, and evaluating $\int_0^{\Tend-t_l}\branch\decay e^{-\decay v}\dd
v = \branch(1-e^{-\decay(\Tend-t_l)})$, gives
\cref{eq:ch05:completedata}. Note the compensator does not depend on $u$:
parents determine only which factor each event's log-intensity lands in.
\end{proof}

The observed-data likelihood \cref{eq:ch05:loglik} is recovered by summing
$e^{\loglik_c}$ over genealogies --- the log-of-sum that destroys
concavity. Augmentation splits it back apart.

\begin{proposition}[E-step: exact posterior parentage]\label{prop:ch05:estep}
Under parameters $(\bcoef,\branch,\decay)$, the posterior over genealogies
given the event times factorizes across events,
$\Prob(u \st t_{1:n}) = \prod_m p_{m,u_m}$, with
\begin{equation}\label{eq:ch05:estep}
p_{m0} \;=\; \frac{\mub(t_m)}{\mub(t_m) + \branch\decay A(m)},
\qquad
p_{ml} \;=\; \frac{\branch\decay\, e^{-\decay(t_m - t_l)}}
{\mub(t_m) + \branch\decay A(m)},
\quad 1 \le l < m,
\end{equation}
with $A(m)$ as in \cref{eq:ch05:Adef} --- so the E-step reuses the Ogata
recursion and costs $O(n)$ beyond the pairwise weights.
\end{proposition}

\begin{proof}
By \cref{lem:ch05:completedata}, $e^{\loglik_c(u)} = C \prod_m w_{m,u_m}$
where $w_{m0} = \mub(t_m)$, $w_{ml} = \branch\decay e^{-\decay(t_m-t_l)}$,
and $C$ collects the compensator factors, which do not involve $u$. Hence
$\Prob(u \st t_{1:n}) \propto \prod_m w_{m,u_m}$, which factorizes with
per-event normalizers $\mub(t_m) + \branch\decay A(m)$. Probabilistically
this is the thinning identity: given $\Ftm$, the event at $t_m$ belongs to
one of the superposed streams with probability proportional to that
stream's intensity.
\end{proof}

\begin{proposition}[M-step]\label{prop:ch05:mstep}
Let $P_0 = \sum_m p_{m0}$ and $P_1 = n - P_0 = \sum_m \sum_{l\ge1} p_{ml}$
be the expected immigrant and offspring counts, and $D = \sum_{m,l\ge1}
p_{ml}(t_m - t_l)$ the expected total parent--child gap. Maximizing
$Q(\cdot) = \E_{u \sim p}[\loglik_c]$ gives, in the interior approximation
that replaces the truncated exposures $1 - e^{-\decay(\Tend-t_l)}$ by $1$:
\begin{equation}\label{eq:ch05:mstep}
\hat\mub = \frac{P_0}{\Tend} \quad (\text{constant baseline}),
\qquad
\hat\branch = \frac{P_1}{n},
\qquad
\hat\decay = \frac{P_1}{D},
\end{equation}
i.e.\ branching = expected children per event, decay = reciprocal of the
$p$-weighted mean parent--child gap. With exact exposures the
$\branch$-update becomes $\hat\branch(\decay) = P_1 \big/ \sum_l
\big(1-e^{-\decay(\Tend-t_l)}\big)$ and $\hat\decay$ solves the scalar
equation $P_1/\decay - D = \hat\branch(\decay)\sum_l (\Tend -
t_l)e^{-\decay(\Tend-t_l)}$.
\end{proposition}

\begin{proof}
$Q$ inherits the separable form of \cref{eq:ch05:completedata} with
indicators replaced by posteriors. The $\branch$-part is $P_1\log\branch -
\branch\sum_l(1-e^{-\decay(\Tend-t_l)})$, concave with the stated root;
the $\decay$-part is $P_1\log\decay - \decay D - \branch\sum_l(1 -
e^{-\decay(\Tend-t_l)})$, whose stationarity condition is the displayed
scalar equation; setting the exposure terms to $1$ (events far from the
window edge) decouples them into \cref{eq:ch05:mstep}. For the
exp-link baseline, the $\bcoef$-part of $Q$ is $\sum_m p_{m0}
\bcoef^\top\xcov(t_m) - \int_0^{\Tend} e^{\bcoef^\top\xcov}$ --- a
weighted Poisson GLM, globally concave. The augmentation thus restores
concavity to the very block that broke it in
\cref{warn:ch05:nonconcave}.
\end{proof}

\begin{theorem}[Monotone ascent]\label{thm:ch05:emascent}
Let $\thetamodel^{(k+1)}$ maximize $Q(\cdot \st \thetamodel^{(k)})$. Then
$\loglik(\thetamodel^{(k+1)}) \ge \loglik(\thetamodel^{(k)})$, with
equality only if $\thetamodel^{(k)}$ already maximizes $Q(\cdot \st
\thetamodel^{(k)})$.
\end{theorem}

\begin{proof}
Write $\pi_{\thetamodel}(u) = \Prob(u \st t_{1:n}; \thetamodel)$ for the
posterior of \cref{prop:ch05:estep}. For any $\thetamodel'$,
$\loglik(\thetamodel') = \E_{\pi_{\thetamodel}}[\loglik_c(\thetamodel';
u)] - \E_{\pi_{\thetamodel}}[\log \pi_{\thetamodel'}(u)]$, since
$\loglik_c(\thetamodel';u) = \loglik(\thetamodel') + \log
\pi_{\thetamodel'}(u)$ holds for every $u$ (Bayes) and the expectation of
a constant is itself. Subtracting the same identity at $\thetamodel$,
\[
\loglik(\thetamodel') - \loglik(\thetamodel)
=
\big[Q(\thetamodel' \st \thetamodel) - Q(\thetamodel \st \thetamodel)\big]
+ \KL{\pi_{\thetamodel}}{\pi_{\thetamodel'}}
\;\ge\;
Q(\thetamodel' \st \thetamodel) - Q(\thetamodel \st \thetamodel)
\;\ge\; 0
\]
for $\thetamodel' = \thetamodel^{(k+1)}$, using nonnegativity of
Kullback--Leibler divergence (Jensen) and the definition of the M-step.
Equality forces both terms to vanish.
\end{proof}

Three practical points, each load-bearing. \emph{Boundary corrections.}
The exact-exposure updates of \cref{prop:ch05:mstep} are not optional
refinements: near the right window edge an event's offspring are
partially unobserved, and the interior approximation deflates $\branch$
while the orphaned early-window events (parents before time $0$) inflate
$\mub$. The right edge is corrected \emph{exactly} inside the M-step ---
the offspring process of an observed parent really is Poisson observed on
$(t_l,\Tend]$ --- while the left edge is only mitigated, by discarding a
burn-in of several kernel half-lives from the immigrant accounting
\citep{veenschoenberg2008}. \emph{The E-step is the endogeneity
diagnostic.} The honest answer to ``how much of deal flow is contagion''
is not the plug-in $\hat\branch$ of \cref{cor:ch05:endogenous} but the
posterior expected offspring share $\frac1n\sum_m (1 - p_{m0})$, together
with the distribution of attributed parent--child gaps: genuine contagion
concentrates attributions at the kernel timescale, while a latent common
factor inflates $\hat\branch$ yet produces attributions with no timescale
structure (\cref{ch05:sec:frailty}). \emph{Constraints are not
automatic.} EM keeps every $\branch \ge 0$ by construction, but nothing
in the parent accounting bounds $\specrad{\hat\Gmat}$ below one --- a rare
source sector attributed many children can push a column past criticality
--- so the subcriticality certificate must be checked and enforced after
each fit, exactly as the discrete-face fitter does. In the multivariate
model the updates read $\hat\Gmat_{sr} = \sum_{m: s_m=s}\sum_{l: s_l=r}
p_{ml} \big/ \sum_{l: s_l = r}\big(1 - e^{-\decay(\Tend - t_l)}\big)$,
expected type-$s$ children per unit of type-$r$ exposure. If profile
likelihoods over the $\decay$-grid show the exponential shape itself
misfits (time-rescaling diagnostics of \cref{ch:gof}), the escalation is
the histogram-kernel nonparametric EM of \citet{lewismohler2011} --- same
E-step, kernel M-step replaced by a binned estimator.

\begin{protocol}[Estimation ladder for the continuous-time sector
layer]\label{prot:ch05:estimation}
(a) Fix a grid of decays $\decay$ spanning well-separated half-lives from
weeks to quarters; do not free $\decay$ jointly with $\branch$ except on
the final reported grid point. (b) At each grid point, fit by EM
(\cref{prop:ch05:estep,prop:ch05:mstep}) with right-edge exposure
corrections and a burn-in of at least four half-lives, multistarted;
cross-check with the direct MLE of
\cref{ch05:sec:likelihood}. (c) Report the profile likelihood over
$\decay$ with its flatness (the ridge diagnostic of \cref{ch:gof}), never
a single $(\hat\branch,\hat\decay)$ point alone. (d) Report
$\specrad{\hat\Gmat}$, the posterior endogenous share, and the
parent--child gap histogram; enforce $\specrad{\hat\Gmat} \le \rhomax =
0.95$ as in the discrete face. (e) All branching estimates carry the
upper-bound language of \cref{prot:ch05:upperbound}.
\end{protocol}

\section{Simulation}\label{ch05:sec:simulation}

Simulation is not an afterthought: multi-step forecasts require it
(history covariates are \emph{internal}, so scenario propagation must
draw events and update features jointly \citep{kalbfleisch2002}), and the
goodness-of-fit machinery of \cref{ch:gof} consumes simulated envelopes.
Three routes, all exact in law.

\begin{lemma}[Thinning bound for exponential kernels]\label{lem:ch05:thinning}
On any interval containing no events and no baseline increase, $t \mapsto
\lam_s(t)$ of \cref{eq:ch05:multi} is nonincreasing. Hence
$\lam^{*} = \sum_s \lam_s(\tau^{+})$ evaluated just after the most recent
event or baseline breakpoint $\tau$ dominates the total intensity until
the next event or breakpoint.
\end{lemma}

\begin{proof}
The excitation part is a nonnegative combination of $e^{-\decay_{sr}(t -
t_m)}$, each strictly decreasing in $t$; the baseline is constant on the
interval by hypothesis. A sum of nonincreasing functions is nonincreasing.
\end{proof}

\emph{Ogata thinning} \citep{lewisshedler1979,ogata1981} then
specializes cleanly: from the current time propose a candidate gap
$\sim\mathrm{Exp}(\lam^{*})$, accept with probability
$\sum_s\lam_s(t_{\mathrm{cand}})/\lam^{*}$, attribute the accepted event
to sector $s$ with probability proportional to $\lam_s$, refresh the bound
at every event \emph{and} at every covariate breakpoint (where the
piecewise-constant baseline may jump up). The Markov embedding keeps the
per-event cost at $O(\abs{\sectors}^2)$ state updates;
\modrepo{eventtime/simulate.py} implements exactly this, including the
breakpoint refresh. \emph{Exact inversion} \citep{dassioszhao2013}
removes rejection entirely for the exponential kernel: $(t, \lam(t))$ is
a piecewise-deterministic Markov process, and the next inter-arrival time
is the minimum of a baseline-driven $\mathrm{Exp}(\bar\mub)$ draw and an
excitation-driven time whose survival function
$\exp\!\big({-}\tfrac{Y}{\decay}(1 - e^{-\decay u})\big)$, $Y$ the current
excess intensity, inverts in closed form (defective with probability
$e^{-Y/\decay}$: the cluster may die). This is the preferred engine for
the long forecast horizons where thinning's bound refreshes dominate cost.
\emph{Cluster simulation} (\cref{thm:ch05:cluster}: draw immigrants, then
recursive Poisson offspring) is the third route; it is the one that
exposes the genealogy, which makes it the reference implementation for
validating EM attributions and edge corrections on synthetic data. All
three must agree in law; the test suite compares them distributionally.

\section{Inhibition and excitation via engineered
covariates}\label{ch05:sec:inhibition}

The market has a signed dynamic: sector-level deal flow is
\emph{self-exciting} (capital attention, momentum), while a firm's own
recent raise is \emph{inhibitory} (a funded firm leaves the market for
months). Regime~B of the rehearsal world builds in exactly this
structure, with a 30-day inhibition half-life against a 7-day contagion
half-life (\cref{ch02:sec:synthetic}). The linear Hawkes family of
\cref{asm:ch05:model} cannot represent the inhibitory half, and the proof
is one paragraph.

\begin{proposition}[Monotonicity: nonnegative kernels cannot
inhibit]\label{prop:ch05:monotone}
Let $\lam(t; \Hist) = \mub(t) + \sum_{t_m \in \Hist,\, t_m < t}
\kernel(t - t_m)$ with $\kernel \ge 0$. Then for every history $\Hist$,
every additional event time $t_0 < t$, and every $t$:
$\lam(t; \Hist \cup \{t_0\}) \ge \lam(t; \Hist)$, and $\lam(t;\Hist) \ge
\mub(t)$. Consequently no parameter choice makes the intensity after an
added event \emph{lower} than before it, nor ever lower than the
baseline: a post-deal decrease in a firm's own intensity is inexpressible
in any linear Hawkes model with nonnegative kernels, univariate or
multivariate, for any $(\branch, \decay)$.
\end{proposition}

\begin{proof}
The intensity is additive in history: $\lam(t; \Hist\cup\{t_0\}) =
\lam(t;\Hist) + \kernel(t - t_0) \ge \lam(t;\Hist)$ since $\kernel \ge
0$; dropping the whole excitation sum gives $\lam \ge \mub$. The infimum
over parameters of the added term is $0$ (attained as $\branch \to 0$),
so the family's expressive range for the effect of an own event is
$[0, \infty)$: ``no effect'' is its best imitation of inhibition. The
positivity transform $\branch = e^{a}$ of
\cref{ch05:sec:likelihood} makes the restriction syntactically invisible
--- the optimizer happily drives $a$ very negative --- which is why this
proposition must be stated rather than discovered in production.
\end{proof}

Three legal routes remain. Each is developed with its model, validity
argument, estimation story, and stability discussion; the decision box at
the end commits.

\subsection*{Route (a): log-link self-modulated intensity (``Hawkes
without Hawkes'')}

Move history inside the link:
\begin{equation}\label{eq:ch05:loglink}
\lam(t) \;=\; \exp\!\big(\bcoef^\top \xcov(t) + \gamma^\top r(t)\big),
\end{equation}
where $r(t)$ is a vector of \emph{engineered history covariates}:
recency indicators and ages (a cooldown flag $\ind{t - \tau_i(t) \le W}$
with $\tau_i(t)$ the last own-deal time, the gap transform $\log(1 +
t-\tau_i(t))$), and decayed event counts
\begin{equation}\label{eq:ch05:decayedcount}
H_{\decay}(t) \;=\; \sum_{t_m < t} e^{-\decay (t - t_m)},
\qquad \dd H_{\decay} = -\decay H_{\decay}\dd t + \dd N(t),
\end{equation}
at a small grid of \emph{fixed} (or profiled) decays, per source sector in
the multivariate case. Negative components of $\gamma$ are inhibition,
positive are excitation: the sign restriction of
\cref{prop:ch05:monotone} evaporates because positivity is carried by the
exponential, not by the kernel. This is the point-process GLM of
\citet{truccolo2005}, and the discrete-time weekly face of it is the
repository's current sector model.

\begin{proposition}[Validity and concavity of the log-link
model]\label{prop:ch05:loglinkvalid}
Let each component of $r$ be a left-continuous functional of
$\Hist_{t^-}$ as above. Then (i) $\lam$ in \cref{eq:ch05:loglink} is
predictable, strictly positive, and a.s.\ finite on the observation
window, so the likelihood \cref{eq:ch02:jacod} is well defined with no
positivity constraints on $(\bcoef,\gamma)$; (ii) $(H_{\decay}, t -
\tau_i(t))$ is a finite-dimensional Markov state, so \cref{eq:ch05:loglink}
depends on history only through a computable summary; (iii) for fixed
decays and windows, the log-likelihood is globally concave in
$(\bcoef,\gamma)$.
\end{proposition}

\begin{proof}
(i) Each feature is adapted and left-continuous (built from events
strictly before $t$), hence predictable; $\exp$ makes $\lam$ positive;
$H_{\decay}(t) \le N(t^-) < \infty$ a.s.\ keeps it finite on the window.
(ii) is the ODE in \cref{eq:ch05:decayedcount} plus the trivial age
dynamics. (iii) The likelihood is $\sum_m (\bcoef^\top\xcov +
\gamma^\top r)(t_m) - \int_0^{\Tend} \exp(\bcoef^\top\xcov +
\gamma^\top r)(u)\dd u$: the event sum is linear in the parameters, and
the compensator is an integral of $\exp(\text{affine})$, convex in the
parameters for every $u$, so its negative is concave; a sum of concave
functions is concave. Estimation is a GLM: the entire fragility of Hawkes
fitting lived in the decay, and fixing it returns the problem to convex
territory. Note what (i) does \emph{not} claim: global-in-time stability
of the process. That is \cref{warn:ch05:mechB}.
\end{proof}

The log-link model is not a retreat from Hawkes; to first order it
\emph{contains} it.

\begin{proposition}[First-order nesting of linear
Hawkes]\label{prop:ch05:nesting}
Fix $\decay$ and constant baseline $\mub$, and write $x(t) =
\branch\decay H_{\decay}(t)/\mub \ge 0$ for the excitation load. The
linear Hawkes intensity $\lam_{\mathrm{lin}} = \mub(1 + x)$ and the
log-link intensity $\lam_{\mathrm{log}} = \mub\, e^{x}$ (i.e.\
\cref{eq:ch05:loglink} with feature $H_\decay$ and coefficient $\gamma =
\branch\decay/\mub$) satisfy
\begin{equation}\label{eq:ch05:nesting}
0 \;\le\; \log \lam_{\mathrm{log}}(t) - \log \lam_{\mathrm{lin}}(t)
\;=\; x(t) - \log\big(1 + x(t)\big) \;\le\; \tfrac{1}{2}\,x(t)^2 .
\end{equation}
Thus a decayed-count feature with $\gamma > 0$ reproduces linear-Hawkes
excitation to first order in the load, with a second-order overshoot; and
$\gamma < 0$ delivers the mirrored soft inhibition that
\cref{prop:ch05:monotone} forbids the linear family. The GLM route nests
soft versions of both.
\end{proposition}

\begin{proof}
$\log\lam_{\mathrm{lin}} = \log\mub + \log(1+x)$ and
$\log\lam_{\mathrm{log}} = \log\mub + x$; the elementary bounds $x -
\tfrac{x^2}{2} \le \log(1+x) \le x$ for $x \ge 0$ give
\cref{eq:ch05:nesting}. In the fitted regime of this repository the load
is small --- covariates dominate and excitation contributes decimal
points of held-out likelihood (\cref{app:results}) --- so the surrogate
error is second-order in an already-small quantity. The overshoot has a
sign, however: the log-link \emph{amplifies} large excursions that the
linear model would meet additively, and that is precisely the seed of the
instability below.
\end{proof}

\begin{warning}[$\specrad{\Gmat}<1$ does not stabilize a log-link count
model (Mechanism B)]\label{warn:ch05:mechB}
For the linear model, subcriticality is a level-free certificate: it
depends only on kernel masses. For log-link feedback it has no such
force. In the weekly discrete face $\log m_{s,t} = a_s + \sum_r B_{sr}
y_{r,t-1}$ (counts $y$ in the exponent), the deterministic skeleton $m
\mapsto \exp(a + Bm)$ has Jacobian $\diag(m^*)B$ at a fixed point $m^*$:
local stability requires $\specrad{\diag(m^*)B} < 1$, an
\emph{operating-point} condition that scales with the level $m^*$. A
stochastic excursion raises the operating point, the local certificate
can silently fail, and the process runs away --- explosion is an
\emph{absorbing rare event}, invisible to in-sample diagnostics. The
rehearsal measured exactly this: a metastable excursion of
$\sim$169{,}000 events in a configuration whose linear-style spectral
radius was nominally far subcritical (\cref{app:results}). The theory
agrees from both sides: geometric ergodicity for log-linear Poisson
autoregressions holds only under contraction conditions on the feedback
coefficients \citep{fokianos2011}, and the continuous-time stability
theory for nonlinear Hawkes requires a Lipschitz link
\citep{bremaud1996} --- $\exp$ is not Lipschitz. Design rules that
follow: history enters the exponent only through \emph{bounded or
saturating} transforms (EWMA levels, $\log(1+\cdot)$ counts, clipped
features), never raw counts; every simulation carries event caps; and the
operating-point diagnostic $\specrad{\diag(\hat m_t)\hat B}$ is monitored
along simulated paths, not certified once at fit time.
\end{warning}

\subsection*{Route (b): rectified additive intensity}

Keep the kernel signed and restore positivity by rectification:
\begin{equation}\label{eq:ch05:rectified}
\lam(t) \;=\; \Big[\, \mub(t)
+ \textstyle\sum_{t_m<t} \kernel^{+}(t-t_m)
- \sum_{t_m<t} \kernel^{-}(t-t_m) \Big]_{+},
\qquad \kernel^{\pm} \ge 0 .
\end{equation}
Validity is immediate ($[\cdot]_+$ of a predictable process is
predictable and nonnegative), and the mathematics is genuinely available:
viewing $\lam = \phi(\mub + \int \kernel_{\mathrm{signed}} \dd N)$ with
$\phi(x)=x_+$ Lipschitz$(1)$, existence, uniqueness, and stability of a
stationary version hold when the total variation of the kernels is
subcritical \citep{bremaud1996}, and the renewal structure of
signed-exponential-kernel processes is worked out in
\citet{costa2020}. Maximum likelihood is feasible for the exponential
kernel because the times at which the rectifier binds (intensity hits
zero, then restarts) are computable in closed form, giving an exact
compensator \citep{bonnet2021}. The costs are structural, not numerical.
The likelihood is nonsmooth in the parameters (kinks wherever a
zero-crossing time moves), estimation loses the GLM concavity of route
(a), and --- decisively for us --- \emph{expectation and rectification do
not commute}: $\E[X]_+ \ne [\E X]_+$ whenever the rectifier binds with
positive probability. The mean-intensity Volterra equation $\xi = \mub +
\kernel * \xi$, whose closed-form exponential-kernel solution is the
engine of the interval-censored MBPP machinery of \cref{ch:censoring},
therefore has no analogue under \cref{eq:ch05:rectified}. Since defect
(C4) forces much of our fitting through interval-censored likelihoods,
route (b) would sever the chapter that keeps timing estimation honest.
It stays on the shelf as the escalation for a genuinely signed
continuous-time ground kernel --- nothing less.

\subsection*{Route (c): move inhibition to the mark factor}

The factorization of \cref{thm:ch03:groundmark} offers the cleanest exit:
keep the \emph{ground} process linear Hawkes over sectors --- positive
momentum only, which is the economically natural sign at the sector level
--- and let firm-level inhibition live in the mark factor, where the
ranker of \cref{ch:ranking} already consumes a cooldown covariate with a
negative coefficient. The firm-level intensity is
\begin{equation}\label{eq:ch05:markfactor}
\lam_i(t) \;=\; \lam_{s(i)}(t)\; p\big(i \st s(i), t, \Hist_{t^-}\big),
\qquad
p(i \st s, t) = \frac{e^{\score_i(t)}}{\sum_{j \in \riskset_{s,t}}
e^{\score_j(t)}},
\quad
\score_i(t) = w^\top \zfeat_i(t) + \eta\, c_i(t),
\end{equation}
with $c_i(t) = \ind{t - \tau_i(t) \le W}$ the cooldown flag and $\eta \le
0$.

\begin{proposition}[Cooldown in the mark factor is firm-level
inhibition]\label{prop:ch05:cooldown}
Under \cref{eq:ch05:markfactor} with $\eta < 0$: (i) $\lam_i$ is
predictable and positive, and the joint likelihood factorizes exactly as
ground $\times$ mark, so both factors are fitted by their own valid
likelihoods; (ii) holding other firms' features fixed, the switch $c_i: 0
\to 1$ multiplies $p(i \st s,t)$ by a factor in $(e^{\eta}, 1)$ --- a
strict decrease --- hence $\lam_i$ strictly decreases while $\lam_s$ is
untouched; (iii) the ground layer retains nonnegative kernels, so the
cluster representation, the $\specrad{\Gmat}$ certificate, and the MBPP
aggregation theory of \cref{ch:censoring} all survive verbatim.
\end{proposition}

\begin{proof}
(i) is \cref{thm:ch03:groundmark} plus predictability of $c_i$ (built
from events strictly before $t$). For (ii), write $p_i(\eta) =
e^{q+\eta}/(D + e^{q+\eta})$ with $q = w^\top\zfeat_i$ and $D =
\sum_{j\ne i} e^{\score_j} > 0$: the map $x \mapsto x/(D+x)$ is strictly
increasing, so $\eta < 0$ strictly decreases $p_i$; the ratio
$p_i(\eta)/p_i(0) = e^{\eta}(D + e^{q})/(D + e^{q+\eta})$ lies strictly
between $e^{\eta}$ and $1$. (iii) is by construction: no signed object
ever touches the ground intensity. What (c) \emph{cannot} express is a
firm's own deal suppressing the \emph{sector aggregate}: the mark factor
only redistributes a fixed sector flow across firms. If aggregate
post-deal suppression is real, it belongs to route (a) at the sector
level, or ultimately to route (b).
\end{proof}

Route (c) is the current production architecture
(\modrepo{sector\_stability.py} + \modrepo{sector\_ranker.py}), and the
rehearsal gives it empirical teeth: the self-inhibition \emph{ordering}
recovered through the ranker cooldown was robust under deliberate
misspecification (rank correlation $0.937$ against regime-B truth,
\cref{app:results}), even though the parametric shape of the cooldown was
wrong.

\begin{decision}{D5.1 --- signed dynamics without signed kernels}
We commit to routes (a) and (c), with division of labor: \emph{(c) owns
firm-level inhibition} --- the linear sector Hawkes ground layer plus the
ranker cooldown covariate, as implemented today --- because it preserves
the cluster representation, the subcriticality certificate, and the MBPP
interval-censored machinery that defect (C4) makes indispensable;
\emph{(a) owns sector-level covariate modulation and timing-feature
models} --- fixed-decay log-link intensities with bounded history
features, as in \modrepo{models/event\_block\_hawkes.py} --- because it
keeps concavity and sign freedom at the price of the operating-point
stability discipline of \cref{warn:ch05:mechB}. Route (b), the rectified
additive model \citep{bonnet2021,costa2020}, is the recorded escalation
if a genuinely signed \emph{continuous-time ground kernel} is ever
required. Reversal triggers: escalate to (b) if the time-rescaling and
PIT diagnostics of \cref{ch:gof} show persistent post-burst
\emph{overprediction of sector aggregates} that route-(a) saturating
features cannot absorb (aggregate suppression, which (c) cannot express
by \cref{prop:ch05:cooldown}), or if the decision layer of
\cref{ch:decision} demands calibrated signed amplification in continuous
time. Restrict (a) --- or shift its mass to (c) --- if the Mechanism-B
operating-point monitor fires in production simulation.
\end{decision}

\section{Frailty versus contagion}\label{ch05:sec:frailty}

A fitted $\hat\branch > 0$ is not proof of contagion. Clustered,
overdispersed event flow has two structural explanations: events beget
events (Hawkes), or a latent common factor modulates everyone's rate
(frailty; the doubly-stochastic alternative). Both produce
overdispersion, autocorrelated counts, and cross-sector comovement, and
from a single realization of moderate length they are close to
observationally equivalent --- the point-process face of the
state-dependence-versus-heterogeneity problem of
\citet{heckmanborjas1980}. The mechanism of confounding is concrete: if a
positively autocorrelated factor $\xi(t)$ multiplies the true baseline
and the model omits it, the score of an excitation loading at zero has
expectation
\begin{equation}\label{eq:ch05:frailtyscore}
\E_0\!\left[\partial_{\gamma}\loglik \big|_{\gamma=0}\right]
\;\propto\;
\int_0^{\Tend} \Cov\big(H_{\decay}(t),\, \xi(t)\big) \dd t \;>\; 0,
\end{equation}
because the decayed count $H_\decay$ is built from past events that
themselves loaded on $\xi$: the fitted excitation absorbs the common
factor. The rehearsal calibrated the damage: under a latent AR($0.90$)
frailty, the freely fitted branching radius came out $2.4\times$ the
truth (\cref{app:results}).

The doctrine that follows is layered. \emph{Covariates first}
\citep{duffie2007}: every observable common driver belongs in
$\mub_s(t)$; comovement explained by covariates was never contagion, and
our ladder runs baseline-first by construction
(\cref{ch:poisson} precedes this chapter). \emph{Latent-factor
escalation} \citep{duffie2009}: with many sectors, frailty and contagion
have different fingerprints --- a common factor loads simultaneously
across sectors, contagion propagates along the event stream at the
kernel's lag structure --- so a latent-factor model (EM or particle
methods over $\xi$) is genuinely estimable; it is an escalation beyond
this chapter that we have deliberately not built. \emph{Both coexist}
\citep{azizpour2018}: the operational conclusion for corporate default
data, after conditioning on covariates, was contagion \emph{and} frailty
--- so the honest posture is not to pick one but to bound. The posterior
parentage diagnostic of \cref{ch05:sec:em} is the sharpest single-sample
probe we have: attributed parent--child gaps concentrated at the kernel
timescale are consistent with contagion; attributions spread evenly in
time indict a common factor.

\begin{protocol}[Branching estimates are upper bounds on
contagion]\label{prot:ch05:upperbound}
Every reported $\hat\branch$, $\hat\Gmat$, or endogenous share is
accompanied by: (i) the covariate set of the baseline it was fitted
against; (ii) the sentence ``upper bound on contagion; unobserved
common factors inflate this estimate'' with the $2.4\times$ AR($0.90$)
calibration of \cref{app:results} cited as the measured severity; (iii)
the parent--child gap histogram of \cref{ch05:sec:em}. No standalone
$\hat\branch$ appears in any deliverable.
\end{protocol}

\section{Defect declaration (D2.1)}\label{ch05:sec:defects}

Against the ledger of \cref{tab:ch02:defects}, for the models of this
chapter (linear sector Hawkes; log-link history models; the mark-factor
cooldown belongs to \cref{ch:ranking}).

\emph{(C1) Administrative censoring: robust (exact).} The likelihood
\cref{eq:ch05:loglik} integrates the compensator to $\Tend$; the open
window contributes its survival factor $\exp(-\int\lam)$ exactly. Unlike
the Poisson case, that factor is \emph{path-dependent} --- it depends on
the realized history through the excitation terms --- but for exponential
kernels it is available in closed form given the history, so nothing is
approximated. Within EM, the right-edge exposure correction of
\cref{prop:ch05:mstep} is the same fact seen from the cluster side: an
observed parent's offspring process really is Poisson observed on
$(t_l,\Tend]$. The subtleties of history-dependent survival under
\emph{interval} censoring are deferred to \cref{ch:censoring}.

\emph{(C2) Left truncation: corrects under assumptions.} All compensator
integrals start at entry times, and risk-set indicators $\atrisk$
multiply intensities as in \cref{eq:ch02:jacod}. The Hawkes-specific
residue is the \emph{left edge}: in-window events whose parents predate
the window are structurally misattributed to immigration, biasing
$\hat\mub$ up and $\hat\Gmat$ down. Burn-in of several half-lives
mitigates; with 90-day-scale kernel mass and a 13.5-year window the
residual bias is small, but short per-sector subwindow fits are not
licensed (\cref{prot:ch05:estimation}).

\emph{(C3) Unlabeled exits: robust at the sector layer, vulnerable at the
firm layer.} Sector counts are complete regardless of zombie firms, so
the ground models of this chapter are immune. Any firm-level intensity
(the route-(a) block model, the ranker) keeps dead firms at risk:
fitted firm hazards are deflated and cooldown effects attenuated, with
the documented sign (downward on firm-level intensities). The
risk-set sensitivity is the largest in the system
(\cref{app:results}); treatments live in \cref{ch:ranking}.

\emph{(C4) Untrusted dates and aggregation: VULNERABLE for $\decay$, with
signed statement.} Date jitter and binning destroy the information that
identifies the timescale: Fisher information in $\decay$ collapses once
bin width exceeds the kernel half-life, while $\branch$ remains estimable
from count balance. The bias is not a small perturbation of $\hat\decay$
but a loss of identification --- the $\branch$--$\decay$ ridge --- and
timing claims from weekly data are vetoed outright. This defect is the
reason \cref{ch:censoring} exists: the interval-censored MBPP likelihood
is the tractable replacement for the exact likelihood this chapter
derived, and \cref{ch:gof} quantifies the bucket-width rule (safe at
$\le 1$ half-life, destroyed at $\ge 4$).

\emph{(C5) Stale firm covariates: pointer.} The sector baselines of this
chapter consume macro and sector-activity covariates, which are never
stale (as-of joined); the engineered history covariates of route (a) are
point-in-time correct by construction. Firm-level covariates entering
$\zfeat_i$ or block-model baselines inherit the attenuation and
informative-observation pathologies of \cref{ch:stale}.

\emph{(C6) Out-of-universe deals: robust at the ground layer, pointer at
the mark layer.} Sector event streams include out-of-universe deals, so
both the excitation sources and the counts this chapter models are
complete. The mark factor must carry an outside option or overstate every
tracked firm's share; that is \cref{ch:universe}'s treatment.

\begin{codehooks}
\emph{Existing.} \modrepo{eventtime/likelihood.py}: exponential-kernel
Jacod likelihood with the per-pair Ogata recursions of
\cref{prop:ch05:recursion} in numerically stable block-cumsum form, plus
analytic gradients. \modrepo{eventtime/estimate.py}: fixed-decay MLE with
the $\alpha = e^{a}$ positivity transform and delta-method standard
errors; note the code parameterizes $\kernel(u) = \alpha e^{-\beta u}$,
so $\branch = \alpha/\beta$, $\decay = \beta$ --- migrating to the
$(\branch,\decay)$ convention of this document is a pending hygiene item.
\modrepo{eventtime/simulate.py}: Ogata thinning specialized via
\cref{lem:ch05:thinning}, with baseline-breakpoint bound refresh.
\modrepo{eventtime/covariates.py}: piecewise-constant exp-link baselines
with closed-form compensators. \modrepo{eventtime/lasso.py}:
FISTA-based $\ell_1$ recovery of sparse excitation matrices
\citep{hansen2015}. \modrepo{sector\_stability.py}: the current discrete
face (weekly log-link count model, nonnegative lag excitation, row-sum
constraint $\rhomax=0.95$, \code{aggregate\_excitation} and
\code{excitation\_spectral\_radius} diagnostics).
\modrepo{models/event\_block\_hawkes.py}: route (a) in production form
--- exp-link event-time intensity with signed recency term ($\rho \ge 0$
on the inhibitory feature) and sector-block excitation, concave by
\cref{prop:ch05:loglinkvalid}.
\emph{To be written.} A continuous-time sector Hawkes module
(\modrepo{eventtime/sector\_hawkes.py}) implementing
\cref{prot:ch05:estimation}: EM fitter with posterior-parentage output
(\cref{prop:ch05:estep}), right-edge exposure corrections, profile-decay
grids, and the $\specrad{\hat\Gmat}$ enforcement; the exact
Dassios--Zhao simulator \citep{dassioszhao2013} alongside thinning; and
the Mechanism-B operating-point monitor of \cref{warn:ch05:mechB} as a
first-class simulation diagnostic.
\end{codehooks}
